\chapter{State Feedback and Optimal Linear Control}

\section{State Feedback and Pole Placement}
State feedback is the cornerstone of modern control theory. Under the assumption that the complete state vector $x(t) \in \R^n$ is accessible for measurement, we can design a linear controller of the form:
\begin{equation}
    u(t) = -K x(t) + r(t) .\label{eq:sf_law}
\end{equation}
where $K \in \R^{m \times n}$ is the state feedback gain matrix and $r(t) \in \R^m$ is the reference input. 

Applying the control law \eqref{eq:sf_law} to the open-loop continuous-time system $\dot{x} = A x + B u$ yields the closed-loop system:
\begin{equation}
    \dot{x}(t) = (A - B K) x(t) + B r(t) .\label{eq:cl_sf}
\end{equation}
Symmetrically, for discrete-time systems $x_{k+1} = A x_k + B u_k$, the closed-loop state equation becomes:
\begin{equation}
    x_{k+1} = (A - B K) x_k + B r_k .\label{eq:cl_sf_dt}
\end{equation}
The primary objective of state feedback is to choose $K$ such that the closed-loop system matrix $A_{cl} = A - B K$ possesses desirable dynamic characteristics, represented by a specified set of eigenvalues (poles).

\begin{mytheorem}{Arbitrary Pole Placement}
The eigenvalues of the closed-loop matrix $A - B K$ can be placed at arbitrary, symmetric locations in the complex plane by selecting an appropriate feedback gain matrix $K$ if and only if the system pair $(A, B)$ is controllable.
\end{mytheorem}

If the system is not controllable, it can be decomposed into controllable and uncontrollable subsystems via the Kalman decomposition:
\begin{equation}
\bar{A} = T A T^{-1} = \begin{bmatrix} A_c & A_{12} \\ 0 & A_{uc} \end{bmatrix}, \quad \bar{B} = T B = \begin{bmatrix} B_c \\ 0 \end{bmatrix} .
\end{equation}
In this coordinate system, the closed-loop matrix becomes:
\begin{equation}
\bar{A} - \bar{B} \bar{K} = \begin{bmatrix} A_c - B_c K_c & A_{12} - B_c K_{uc} \\ 0 & A_{uc} \end{bmatrix} .
\end{equation}
The eigenvalues of the closed-loop system are the union of the eigenvalues of $A_c - B_c K_c$ and the eigenvalues of $A_{uc}$. Since the uncontrollable eigenvalues (eigenvalues of $A_{uc}$) are unaffected by the feedback gain $K$, arbitrary pole placement is impossible. However, the closed-loop system can still be stabilized if and only if all eigenvalues of $A_{uc}$ are stable (Hurwitz in continuous-time or Schur in discrete-time). This weaker property is called \textbf{stabilizability}.

\subsection{Ackermann's Formula}
For single-input single-output (SISO) systems ($u(t) \in \R^1$), a direct closed-form formula for computing the state feedback gain $K$ is known as \textbf{Ackermann's Formula}.

Let the desired closed-loop characteristic polynomial be:
\begin{equation}
\alpha_{cl}(s) = s^n + \alpha_{n-1} s^{n-1} + \dots + \alpha_1 s + \alpha_0 .
\end{equation}
We wish the closed-loop matrix $A_{cl} = A - B K$ to satisfy its own characteristic equation, meaning $\alpha_{cl}(A - B K) = 0$ by the Cayley-Hamilton theorem.

\begin{mytheorem}{Ackermann's Formula}
For a controllable SISO pair $(A, B)$, the unique state feedback gain vector $K \in \R^{1 \times n}$ that places the eigenvalues of $A - B K$ at the roots of $\alpha_{cl}(s) = 0$ is:
\begin{equation}
    K = \begin{bmatrix} 0 & 0 & \dots & 0 & 1 \end{bmatrix} \mathcal{C}^{-1} \alpha_{cl}(A) .\label{eq:ackermann}
\end{equation}
where $\mathcal{C} = \begin{bmatrix} B & AB & A^2B & \dots & A^{n-1}B \end{bmatrix}$ is the controllability matrix, and $\alpha_{cl}(A)$ is the matrix polynomial:
\begin{equation}
\alpha_{cl}(A) = A^n + \alpha_{n-1} A^{n-1} + \dots + \alpha_1 A + \alpha_0 I .
\end{equation}
\end{mytheorem}

\begin{proof}
Let us analyze the powers of the closed-loop matrix $A_{cl} = A - BK$:
\begin{align}
    A_{cl} &= A - BK \nonumber \\
    A_{cl}^2 &= (A - BK)^2 = A^2 - ABK - BKA_{cl} \nonumber \\
    A_{cl}^3 &= A^3 - A^2BK - ABKA_{cl} - BKA_{cl}^2 \nonumber \\
    &\vdots \nonumber \\
A_{cl}^n &= A^n - A^{n-1}BK - A^{n-2}BKA_{cl} - \dots - BKA_{cl}^{n-1} .
\end{align}
Now, evaluate the desired closed-loop characteristic polynomial at the matrix $A$:
\begin{align}
    \alpha_{cl}(A) &= A^n + \alpha_{n-1} A^{n-1} + \dots + \alpha_1 A + \alpha_0 I \nonumber \\
&= A_{cl}^n + \alpha_{n-1} A_{cl}^{n-1} + \dots + \alpha_1 A_{cl} + \alpha_0 I + \sum_{j=1}^n A^{n-j} B K A_{cl}^{j-1} .
\end{align}
By the Cayley-Hamilton theorem, $\alpha_{cl}(A_{cl}) = A_{cl}^n + \alpha_{n-1} A_{cl}^{n-1} + \dots + \alpha_1 A_{cl} + \alpha_0 I = 0$. Therefore, the expression simplifies to:
\begin{equation}
\alpha_{cl}(A) = \sum_{j=1}^n A^{n-j} B K A_{cl}^{j-1} = \begin{bmatrix} B & AB & \dots & A^{n-1}B \end{bmatrix} \begin{bmatrix} K A_{cl}^{n-1} \\ K A_{cl}^{n-2} \\ \vdots \\ K \end{bmatrix} .
\end{equation}
We recognize the controllability matrix $\mathcal{C}$ on the right-hand side. Multiplying by $\mathcal{C}^{-1}$ from the left yields:
\begin{equation}
\mathcal{C}^{-1} \alpha_{cl}(A) = \begin{bmatrix} K A_{cl}^{n-1} \\ K A_{cl}^{n-2} \\ \vdots \\ K \end{bmatrix} .
\end{equation}
To isolate $K$, we extract the bottom row of this matrix equation. Multiplying by the row vector $\begin{bmatrix} 0 & 0 & \dots & 1 \end{bmatrix}$ selects the last row, which is exactly $K$:
\begin{equation}
K = \begin{bmatrix} 0 & 0 & \dots & 0 & 1 \end{bmatrix} \mathcal{C}^{-1} \alpha_{cl}(A) .
\end{equation}
This completes the derivation of Ackermann's formula.
\end{proof}

\section{Continuous-Time Linear Quadratic Regulator (LQR)}
While pole placement allows arbitrary placement of eigenvalues, it does not provide guidance on *where* to place them. Moving poles far into the left-half plane stabilizes the states quickly but demands massive control effort (actuator saturation). Conversely, placing poles close to the imaginary axis saves control energy but results in slow, oscillatory transient response.

The \textbf{Linear Quadratic Regulator (LQR)} resolves this fundamental trade-off by casting controller design as an optimization problem.

Consider the continuous-time LTI system:
\begin{equation}
\dot{x}(t) = A x(t) + B u(t), \quad x(0) = x_0 .
\end{equation}
We seek an optimal control input $u^*(t)$ that minimizes the infinite-horizon quadratic cost functional:
\begin{equation}
    J = \int_0^\infty \left( x(t)^\trp Q x(t) + u(t)^\trp R u(t) \right) dt .\label{eq:lqr_cost}
\end{equation}
where:
- $Q = Q^\trp \ge 0$ is the state weighting matrix, penalizing deviations of the states.
- $R = R^\trp > 0$ is the control weighting matrix, penalizing control effort.

\subsection{Derivation via the Hamiltonian and Variational Calculus}
To minimize $J$ subject to the state dynamic constraint $\dot{x} = Ax + Bu$, we formulate the Hamiltonian function:
\begin{equation}
H(x, u, \lambda) = x^\trp Q x + u^\trp R u + \lambda^\trp (A x + B u) .
\end{equation}
where $\lambda(t) \in \R^n$ is the vector of Lagrange multipliers (or co-states).

According to Pontryagin's Minimum Principle, the optimal state trajectory $x^*(t)$, optimal control $u^*(t)$, and optimal co-state $\lambda^*(t)$ must satisfy the following necessary conditions:
\begin{enumerate}
    \item \textbf{Optimality Condition}:
    \begin{equation}
        \frac{\partial H}{\partial u} = 2 R u + B^\trp \lambda = 0 \implies u^*(t) = -\frac{1}{2} R^{-1} B^\trp \lambda(t) .\label{eq:lqr_opt_u}
    \end{equation}
    \item \textbf{Co-state Equation}:
    \begin{equation}
        \dot{\lambda} = -\frac{\partial H}{\partial x} = -2 Q x - A^\trp \lambda .\label{eq:lqr_costate}
    \end{equation}
    \item \textbf{State Equation}:
    \begin{equation}
        \dot{x} = \frac{\partial H}{\partial \lambda} = A x + B u = A x - \frac{1}{2} B R^{-1} B^\trp \lambda .\label{eq:lqr_state}
    \end{equation}
\end{enumerate}

To solve this two-point boundary value problem, we propose that the co-state is linearly related to the state vector:
\begin{equation}
    \lambda(t) = 2 P(t) x(t) .\label{eq:lqr_sweep}
\end{equation}
where $P(t)$ is a symmetric matrix. For the infinite-horizon problem, the system is time-invariant, meaning $P(t) = P$ is a constant symmetric matrix.

Differentiating \eqref{eq:lqr_sweep} with respect to time yields:
\begin{equation}
\dot{\lambda} = 2 P \dot{x} .
\end{equation}
Substituting the state equation \eqref{eq:lqr_state} and the co-state equation \eqref{eq:lqr_costate} into this relation gives:
\begin{equation}
-2 Q x - A^\trp (2 P x) = 2 P \left( A x - \frac{1}{2} B R^{-1} B^\trp (2 P x) \right) .
\end{equation}
Dividing by 2 and moving all terms to one side:
\begin{equation}
\left( A^\trp P + P A - P B R^{-1} B^\trp P + Q \right) x(t) = 0 .
\end{equation}
For this relation to hold for all possible state trajectories $x(t)$, the matrix term in parentheses must vanish. This yields the continuous-time \textbf{Algebraic Riccati Equation (ARE)}:
\begin{equation}
    A^\trp P + P A - P B R^{-1} B^\trp P + Q = 0 .\label{eq:lqr_are}
\end{equation}
Substituting $\lambda(t) = 2 P x(t)$ back into the optimality condition \eqref{eq:lqr_opt_u} yields the optimal state feedback law:
\begin{equation}
    u^*(t) = -K x(t), \quad K = R^{-1} B^\trp P .\label{eq:lqr_gain}
\end{equation}

\begin{mytheorem}{Infinite-Horizon LQR Stability}
If the system pair $(A, B)$ is stabilizable and the pair $(A, Q^{1/2})$ is detectable, then:
\begin{enumerate}
    \item The Algebraic Riccati Equation \eqref{eq:lqr_are} has a unique symmetric positive definite solution $P \succ 0$.
    \item The optimal closed-loop system matrix $A - B K = A - B R^{-1} B^\trp P$ is Hurwitz (strictly stable).
\end{enumerate}
\end{mytheorem}

\begin{proof}
Let $P \succ 0$ satisfy the ARE \eqref{eq:lqr_are}. We can rewrite the ARE by adding and subtracting $P B R^{-1} B^\trp P$:
\begin{align}
    A^\trp P + P A - P B R^{-1} B^\trp P + Q &= 0 \nonumber \\
    (A - B R^{-1} B^\trp P)^\trp P + P (A - B R^{-1} B^\trp P) + P B R^{-1} B^\trp P + Q &= 0 \nonumber \\
    (A - B K)^\trp P + P (A - B K) &= -(Q + K^\trp R K) .\label{eq:are_rewritten}
\end{align}
Let us define a quadratic Lyapunov function candidate $V(x) = x^\trp P x$. Differentiating $V(x)$ along the closed-loop system trajectories $\dot{x} = (A - B K) x$ yields:
\begin{equation}
\dot{V}(x) = x^\trp \left( (A - B K)^\trp P + P (A - B K) \right) x = -x^\trp (Q + K^\trp R K) x .
\end{equation}
Since $Q \ge 0$ and $R > 0$, the matrix $Q + K^\trp R K \ge 0$, which implies $\dot{V}(x) \le 0$. By Lyapunov's stability theory, the closed-loop system is stable in the sense of Lyapunov. To establish asymptotic stability, we invoke LaSalle's Invariance Principle or detectability of $(A, Q^{1/2})$.
Suppose there exists a non-zero trajectory $x(t)$ such that $\dot{V}(x(t)) = 0$ for all $t$. This requires:
\begin{equation}
x(t)^\trp Q x(t) = 0 \quad \text{and} \quad x(t)^\trp K^\trp R K x(t) = 0 .
\end{equation}
Since $R \succ 0$, $K x(t) = 0$, meaning the control input $u(t) = -K x(t) = 0$. The dynamics reduce to the unforced system $\dot{x} = A x$. Furthermore, $x(t)^\trp Q x(t) = 0$ implies $Q^{1/2} x(t) = 0$.
By detectability of $(A, Q^{1/2})$, any state trajectory that satisfies $Q^{1/2} x(t) = 0$ for all $t$ under $\dot{x} = Ax$ must converge to zero: $\lim_{t\to\infty} x(t) = 0$.
Thus, the only invariant set where $\dot{V}(x) = 0$ is the origin $x=0$, proving asymptotic stability of $A - BK$.
\end{proof}

\section{Discrete-Time Linear Quadratic Regulator}
For discrete-time systems, the optimization problem is formulated over a discrete sum. Let the system be:
\begin{equation}
x_{k+1} = A x_k + B u_k, \quad x_0 \text{ given} .
\end{equation}
We wish to minimize the discrete infinite-horizon cost functional:
\begin{equation}
    J = \sum_{k=0}^\infty \left( x_k^\trp Q x_k + u_k^\trp R u_k \right) .\label{eq:lqr_dt_cost}
\end{equation}
with $Q = Q^\trp \ge 0$ and $R = R^\trp \succ 0$.

\subsection{Derivation via Dynamic Programming and the Bellman Equation}
We derive the discrete optimal controller using Bellman's Principle of Optimality. Let $V_k(x_k)$ represent the optimal cost-to-go from step $k$ to infinity starting at state $x_k$:
\begin{equation}
V_k(x_k) = \min_{u_k, u_{k+1}, \dots} \sum_{j=k}^\infty \left( x_j^\trp Q x_j + u_j^\trp R u_j \right) .
\end{equation}
The Principle of Optimality allows us to split the sum and write the recursive \textbf{Bellman Equation}:
\begin{equation}
    V_k(x_k) = \min_{u_k} \left\{ x_k^\trp Q x_k + u_k^\trp R u_k + V_{k+1}(x_{k+1}) \right\} .\label{eq:bellman}
\end{equation}
We propose that the optimal cost-to-go is quadratic in the state:
\begin{equation}
V_{k+1}(x_{k+1}) = x_{k+1}^\trp P_{k+1} x_{k+1} .
\end{equation}
where $P_{k+1} \succ 0$ is a symmetric matrix. Substituting the system dynamics $x_{k+1} = A x_k + B u_k$ into the Bellman equation \eqref{eq:bellman}:
\begin{equation}
V_k(x_k) = \min_{u_k} \left\{ x_k^\trp Q x_k + u_k^\trp R u_k + (A x_k + B u_k)^\trp P_{k+1} (A x_k + B u_k) \right\} .
\end{equation}
Expanding the terms inside the minimization brackets:
\begin{align}
    \mathcal{J}_k(u_k) &= x_k^\trp Q x_k + u_k^\trp R u_k + (x_k^\trp A^\trp + u_k^\trp B^\trp) P_{k+1} (A x_k + B u_k) \nonumber \\
&= x_k^\trp (Q + A^\trp P_{k+1} A) x_k + 2 u_k^\trp B^\trp P_{k+1} A x_k + u_k^\trp (R + B^\trp P_{k+1} B) u_k .
\end{align}
Since $R \succ 0$ and $P_{k+1} \succ 0$, the matrix $R + B^\trp P_{k+1} B$ is strictly positive definite. The function $\mathcal{J}_k(u_k)$ is strictly convex in $u_k$. The global minimum is found by setting the derivative with respect to $u_k$ to zero:
\begin{equation}
\frac{\partial \mathcal{J}_k}{\partial u_k} = 2 (R + B^\trp P_{k+1} B) u_k + 2 B^\trp P_{k+1} A x_k = 0 .
\end{equation}
Solving for $u_k$ yields the optimal control law:
\begin{equation}
    u_k^* = -K_k x_k, \quad K_k = (R + B^\trp P_{k+1} B)^{-1} B^\trp P_{k+1} A .\label{eq:lqr_dt_gain}
\end{equation}
Now, substitute $u_k^*$ back into the cost function to determine the updated cost-to-go $V_k(x_k) = x_k^\trp P_k x_k$:
\begin{align}
    x_k^\trp P_k x_k &= x_k^\trp (Q + A^\trp P_{k+1} A) x_k - 2 x_k^\trp A^\trp P_{k+1} B (R + B^\trp P_{k+1} B)^{-1} B^\trp P_{k+1} A x_k \nonumber \\
    &\quad + x_k^\trp A^\trp P_{k+1} B (R + B^\trp P_{k+1} B)^{-1} (R + B^\trp P_{k+1} B) (R + B^\trp P_{k+1} B)^{-1} B^\trp P_{k+1} A x_k \nonumber \\
&= x_k^\trp \left( A^\trp P_{k+1} A - A^\trp P_{k+1} B (R + B^\trp P_{k+1} B)^{-1} B^\trp P_{k+1} A + Q \right) x_k .
\end{align}
For this to hold for all $x_k$, we obtain the backward difference Riccati equation:
\begin{equation}
P_k = A^\trp P_{k+1} A - A^\trp P_{k+1} B (R + B^\trp P_{k+1} B)^{-1} B^\trp P_{k+1} A + Q .
\end{equation}
For the infinite-horizon problem, the matrix converges to a steady-state value $\lim_{k\to-\infty} P_k = P$, yielding the \textbf{Discrete Algebraic Riccati Equation (DARE)}:
\begin{equation}
    P = A^\trp P A - A^\trp P B (R + B^\trp P B)^{-1} B^\trp P A + Q .\label{eq:lqr_dare}
\end{equation}
The optimal steady-state state feedback gain is:
\begin{equation}
    K = (R + B^\trp P B)^{-1} B^\trp P A .\label{eq:lqr_dt_optimal_gain}
\end{equation}

Similar to the continuous-time case, if $(A, B)$ is stabilizable and $(A, Q^{1/2})$ is detectable, the DARE \eqref{eq:lqr_are} has a unique positive definite solution $P \succ 0$, and the closed-loop matrix $A - B K$ is strictly Schur (stable).

\section{Robust State Feedback Synthesis via LMIs}
When the system matrices $A$ and $B$ are subject to polytopic uncertainty:
\begin{equation}
[A(t) \quad B(t)] \in \text{Co}\{[A_1 \quad B_1], \dots, [A_M \quad B_M]\} .
\end{equation}
the closed-loop matrix under state feedback $u = -K x$ is $A(t) - B(t)K$. For robust quadratic stability, we seek a single symmetric positive definite matrix $P \succ 0$ and a single feedback gain $K$ such that:
\begin{equation}
    (A_i - B_i K)^\trp P + P (A_i - B_i K) \prec 0, \quad \forall i = 1, \dots, M .\label{eq:robust_sf_bmi}
\end{equation}
The inequality \eqref{eq:robust_sf_bmi} is a \textbf{Bilinear Matrix Inequality (BMI)} because it contains bilinear terms involving the products of the decision variables $P$ and $K$ (e.g., $P B_i K$). BMIs are non-convex and notoriously difficult to solve numerically.

To transform this BMI into a solvable \textbf{Linear Matrix Inequality (LMI)}, we perform a congruence transformation and change of variables:
\begin{enumerate}
    \item Define $W = P^{-1} \succ 0$.
    \item Pre- and post-multiply the inequality \eqref{eq:robust_sf_bmi} by $W$:
    \begin{equation}
W (A_i - B_i K)^\trp P W + W P (A_i - B_i K) W \prec 0 .
    \end{equation}
    Since $WP = I$, this simplifies to:
    \begin{equation}
W A_i^\trp - W K^\trp B_i^\top + A_i W - B_i K W \prec 0 .
    \end{equation}
    \item Define the new auxiliary variable:
    \begin{equation}
Y = K W \in \R^{m \times n} \implies Y^\trp = W K^\trp .
    \end{equation}
    Substituting $Y$ and $Y^\trp$ into the inequality yields:
    \begin{equation}
A_i W + W A_i^\trp - B_i Y - Y^\trp B_i^\trp \prec 0 .
    \end{equation}
\end{enumerate}

This leads directly to the robust state-feedback synthesis theorem.

\begin{mytheorem}{Robust LMI State Feedback Synthesis}
The polytopic uncertain system is robustly quadratically stabilizable via a constant state feedback gain $K$ if there exist a symmetric matrix $W \in \R^{n \times n}$ and a matrix $Y \in \R^{m \times n}$ satisfying the LMI system:
\begin{align}
    W &\succ 0 \label{eq:lmi_w} \\
    A_i W + W A_i^\trp - B_i Y - Y^\trp B_i^\trp &\prec 0, \quad \forall i = 1, \dots, M .\label{eq:lmi_vertex}
\end{align}
Once a solution pair $(W, Y)$ is found, the robust state feedback gain is reconstructed as:
\begin{equation}
    K = Y W^{-1} .\label{eq:lmi_k_reconstruct}
\end{equation}
\end{mytheorem}
This transformation is a cornerstone of robust control. It shifts the numerical complexity from non-convex bilinear optimizations to highly efficient, convex semi-definite programming (SDP), which can be solved globally in polynomial time.
